\documentclass[11pt,letterpaper]{article}
\usepackage[T1]{fontenc}
\usepackage{lmodern,amsmath,amssymb,amsthm,booktabs,array,geometry,fancyhdr,xcolor,hyperref}
\geometry{top=.88in,bottom=.86in,left=.94in,right=.94in}
\definecolor{ink}{HTML}{18344A}
\definecolor{linkcol}{HTML}{1D6172}
\hypersetup{colorlinks=true,urlcolor=linkcol,linkcolor=linkcol,citecolor=linkcol,
pdftitle={Sparse equations for singular moduli and the sign of weighted prime sums}}
\urlstyle{same}
\setlength{\headheight}{14pt}
\pagestyle{fancy}\fancyhf{}
\fancyhead[L]{\small\color{ink}Sparse CM equations and weighted prime sums}
\fancyhead[R]{\small\color{ink}Research note}
\fancyfoot[C]{\small\thepage}
\setlength{\emergencystretch}{2em}\raggedbottom
\newtheorem{theorem}{Theorem}[section]
\newtheorem{lemma}[theorem]{Lemma}
\newtheorem{proposition}[theorem]{Proposition}
\newcommand{\upfile}[2]{\href{https://github.com/lsernau/grh-lean/blob/72fead102f663d5c7a407964a13dc0cf9d977207/#1}{#2}}
\begin{document}
\thispagestyle{plain}
\noindent{\small\color{linkcol}\scshape Analytic and algebraic number theory}\hfill{\small November 25, 2026}\par
\vspace{.7em}
\noindent{\fontsize{24}{28}\selectfont\bfseries\color{ink}Sparse equations for singular moduli\\ and the sign of weighted prime sums\par}
\vspace{.6em}
\noindent{\large Applications and extensions of the Dirichlet GRH development}\par
\vspace{.7em}
\begin{abstract}
A singular modulus of degree at least four cannot satisfy a rational polynomial with four or fewer nonzero coefficients, of any degree. We prove this using the behavior of its conjugates at 7 and 13, elementary determinants, and an exhaustive finite certificate. The proof also closes all 79 residual cases in an earlier ProveIt report. A single-prime argument completes the published unconditional reductions of Riffaut's trinomial conjecture. The minimum support of a relation grows at least as $|\Delta|^{1/4-\epsilon}$ unconditionally, and at least as the square root of the class number using Dirichlet GRH.

For two questions of Suzuki about prime powers, we prove a fixed-modulus sign reversal between $Q=138$ and $139$ and a lower logarithmic density exceeding $99.677\%$ for a negative unsmoothed error. For general fixed $Q\ge139$ the eventual sign in his equation (29) detects precisely a central Dirichlet zero of conductor at most $Q$. A separate application of localized averaging proves simultaneous convex balance of prime-power errors in every sufficiently large multiplicative interval. We distinguish the Lean theorems, published analytic inputs, and the new paper arguments with their finite checks.
\end{abstract}

\section{Which parts of the Lean proof transfer?}
The \href{https://github.com/lsernau/grh-lean/tree/72fead102f663d5c7a407964a13dc0cf9d977207}{examined revision} proves Mathlib's Riemann hypothesis and zero location for all Dirichlet characters over $\mathbb Q$. The readable edition's \upfile{grh/src/Statements.lean}{\texttt{Statements.lean}} imports Mathlib alone. Its \upfile{grh/src/Main.lean}{\texttt{Main.lean}} exports five results: ordinary RH, its explicit zeta-zero formulation, half-plane nonvanishing for every modulus $N\ge1$ and character $\chi$ (with the principal pole excepted), zero location in $0<\Re s<1$, and the primitive-character specialization. Central vanishing is allowed. The \upfile{grh/README.md}{upstream verification record} reports Lean 4.35.0, Mathlib \texttt{56a00585}, and the three axioms \texttt{propext}, \texttt{Classical.choice}, \texttt{Quot.sound}. The authors report kernel and independent-checker verification. Our inspection covers declarations and source identities. We have not independently rebuilt either complete Lean project or formalized the extensions in this paper.

\paragraph{The arithmetic estimate.}
In the notation of the \upfile{grh/GUIDE.md}{reading guide}, the proof compares two estimates for the same twisted lattice of automorphic forms. Short integral vectors from theta constructions and a spanning result from a second moment give one inequality. Determinants of chains of theta lifts give another. Their cancellation in \upfile{grh/src/New/Reduction/MainInequality.lean}{\texttt{star\_inequality}} bounds a remainder in
\[
 \sum_{\substack{n\le x\\n\equiv1\pmod q}}\Lambda(n)(1-\chi_K(n))
            (1/n-1/x)-\frac{\log x}{\varphi(q)}-c_q.
\]
The choice of a chain of length comparable with $\sqrt x$ supplies a saving of its reciprocal. The \upfile{grh/src/New/Reduction/OneSidedBound.lean}{one-sided estimate} is combined with auxiliary quadratic twists and a Landau argument to exclude a putative zero. The twist matters because residues involving different $L$-functions can cancel. The selection of $K$, the theta data and the denominator estimates are family-specific. For a general nonabelian Artin or modular $L$-function, a suitable inequality with its Frobenius or Hecke coefficients would be an additional arithmetic theorem. We obtain no such zero-location assertion here.

\paragraph{Two paths out of this argument.}
First, observing what happens at each nonsplit prime is also fruitful for singular moduli. For this purpose a single prime where several conjugates have the same reduction can be decisive. This comes from CM reduction theory; the local matrix below is a generalized Vandermonde, distinct from the theta-lift determinant used upstream.
Second, the \upfile{grh-compact/src/New/Endgame/OneSidedCriterion.lean}{compact edition's endgame} explicitly points to Suzuki's study of signs. There the Mellin transform turns zero multiplicity into a pole residue, and high Laplace moments localize a positive average in a long window. We use the residue at the center to compute a sign, and generalize localization to vector-valued errors. These uses have different dependencies: the CM classification survives removal of the analytic input, while the sign and convex-hull theorems use its zero-location conclusion.

\paragraph{Starting point in ProveIt.}
We compared the \href{https://github.com/VladimirReshetnikov/ProveIt/tree/c59bde1a0f52eee7a7a8dc6b326072f9b71f76a6/docs/incoming}{incoming reports at \texttt{c59bde1a}}. They investigate, among other topics, CM newform coefficient fields, abelian-field classifications, nonabelian one-sided zero tests, and several smoothed prime races. The direct predecessor for Sections~2--3 is \href{https://github.com/VladimirReshetnikov/ProveIt/blob/c59bde1a0f52eee7a7a8dc6b326072f9b71f76a6/docs/incoming/grh_cm_prime_applications.zip}{\texttt{grh\_cm\_prime\_applications.zip}}. It supplied the residual list, the archimedean separation idea and the prior asymptotic benchmark. In the present proof the exhaustive reduction, disk certificates, and growing-support arguments can be checked independently of that archive. We make no claim to bibliographic priority for the new formulations.
\section{Four-term equations for singular moduli}
For an imaginary quadratic order of discriminant $\Delta=-D$, let $h=h(\Delta)$ denote its proper ideal class number, and let $s(\Delta)$ be the minimum number of nonzero coefficients of a nonzero rational polynomial vanishing at a singular modulus of that discriminant. These moduli are the $h$ distinct conjugates of a single algebraic integer. Thus $s$ is independent of the chosen conjugate and $s\le h+1$. A sparse multiple of the minimal polynomial can have arbitrarily large degree a priori.

\begin{theorem}\label{th:four}
For every imaginary quadratic order of class number at least four,
$s(\Delta)\ge5$. In particular, $s(\Delta)=5$ in degree four.
Equivalently, a singular modulus admits a rational annihilator of support at most four exactly when its degree is at most three. The same conclusions hold for polynomial coefficients in $K=\mathbb Q(\sqrt\Delta)$.
\end{theorem}

This answers the first further research question in the incoming report \emph{Sparse CM equations, prime-generating constants, and additive representations}, whose remaining list comprised 50 discriminants of class number four and 29 of class number five. We give an independent reduction to finite calculations, so the conclusion requires neither the earlier large computation nor its GRH cutoff. The equality over $K$ follows from the main theorem of complex multiplication: $[K(j):K]=[\mathbb Q(j):\mathbb Q]=h$. The powers of $j$ satisfying any specified linear dependence over $K$ therefore also satisfy one over $\mathbb Q$ on a subset of the same support.

\subsection{A determinant test inside a \texorpdfstring{$p$}{p}-adic disk}
Normalize $v=v_p$ by $v(p)=1$, and work in $\mathbb C_p$. An open disk around a nonzero center $c$ is measured in relative coordinates $v(z/c-1)$. The following proof makes the thresholds required at 2 and 3 explicit. General $p$-adic root bounds for sparse polynomials were developed by Lenstra and by Avenda\~no and Krick \cite{lenstra,avk}. Our application uses their setting of a single residue disk.

\begin{lemma}\label{lem:disk}
Let $k\ge2$ and choose $r\ge 1/(p-1)$ such that for all $0\le j<k$, $d\ge1$,
\[
             v_p((j+d)!/j!)\le dr.
\]
If $k$ distinct numbers $z_i$ all satisfy $v(z_i/c-1)>r$, then
$\det(z_i^{n_j})\ne0$ for every choice of distinct nonnegative integers $n_j$.
For $p\ge k$, $r=1/(p-k+1)$ is valid. For $k=4$ or $5$, one may use $r=2$ at $p=2$ and $r=1$ at $p=3$.
\end{lemma}
\begin{proof}
After column scaling take $z_i/c=\exp(y_i)$, where the $p$-adic logarithms $y_i$ are distinct and have valuation greater than $r$. The logarithm and exponential are inverse isometries on this disk. Choose $r'>r$ below the least finite $v(y_i)$.
Expand each exponential in the matrix $\exp(y_i n_j)$. The finite Cauchy--Binet formula followed by a limit gives terms indexed by $0\le b_0<\cdots<b_{k-1}$,
\[
 \frac{\det(y_i^{b_j})\det(n_i^{b_j})}{\prod_{j=0}^{k-1} b_j!}.
\]
The term with $b_j=j$ is nonzero. Divide by it. Each alternant divided by its ordinary Vandermonde is an integer polynomial, homogeneous of degree $m=\sum_j(b_j-j)$. Indeed each difference of two variables divides the alternant in $\mathbb Z[X_1,\ldots,X_k]$ and these nonassociate irreducibles have product the Vandermonde. At the $n_i$ the quotient is integral, and at the $y_i$ it has valuation at least $mr'$. The quotient of factorials costs at most $mr$ by hypothesis. All nonleading ratios thus have valuation at least $m(r'-r)>0$, tending to infinity with $m$. The leading unit survives the sum.

For $p\ge k$ the factorial cost is zero when $j+d<p$. Otherwise, by Legendre's digit-sum formula, it is at most $(j+d-1)/(p-1)\le d/(p-j)$ for $j\ge1$ and $d\ge p-j$. The case $j=0$ follows directly from $v_p(d!)\le d/(p-1)$. This gives the advertised bound. For $j\le4$ and $d\ge3$ the same formula gives $(j+d-1)/(p-1)\le 2d$ at 2 and $\le d$ at 3. For $d=1,2$ check the factors $j+1,j+2\le6$ directly.
\end{proof}

To exclude polynomials with fewer than $k$ coefficients, append unused exponents to make $k$ columns. The assertion therefore controls all degrees and coefficients.

\subsection{A one-prime proof of Riffaut's conjecture}
\begin{proposition}\label{prop:three}
A singular modulus of degree at least three satisfies no nonzero rational polynomial with at most three nonzero coefficients. This conclusion has a proof independent of GRH.
\end{proposition}
\begin{proof}
The unconditional Theorems 1.3 and 1.4 of Bilu, Luca and Pizarro-Madariaga say that a trinomial discriminant must be fundamental and satisfy $D>10^{11}$. Their Corollary 5.2 gives $p>4\sqrt D/\log D$ for every split prime $p$ in its quadratic field \cite{blp}. Hence 7 is inert or ramified. A binomial is included in their definition by allowing a zero middle coefficient after normalization, and alternatively one can square the normalized binomial to obtain a trinomial. The roots in question are nonzero.

Fix a prime above 7 in the Hilbert class field. By Deuring's reduction theorem all the conjugate $j$-values reduce to supersingular invariants \cite{eoy}. All are algebraic integers with potentially good reduction. The Hilbert class field is unramified over $K$, so its absolute ramification index at 7 is at most two. The only supersingular $j$ in characteristic 7 is $1728\equiv6$. Thus each conjugate has $v_7(j/6-1)\ge1/2$. Three distinct conjugates in this disk contradict Lemma~\ref{lem:disk} with $k=3$ and $r=1/5$.
\end{proof}
The conditional inference from Dirichlet GRH is the published Theorem 1.2 of \cite{blp}. The local last step above uses the same published restrictions on a hypothetical counterexample and needs no theorem about zeros.

\subsection{Split primes, nonsplit primes and the conductor}
Write $\Delta=f^2 d_K$ with $d_K$ fundamental. The ring class field of this order is unramified over $K$ away from $f$ \cite{cox}. At $p=7,13$ the supersingular invariant is unique and nonzero. To check the two values directly, in characteristic $p>3$ use the Hasse invariant of $y^2=x^3+Ax+B$, the coefficient of $x^{p-1}$ in $(x^3+Ax+B)^{(p-1)/2}$. In characteristic 7 it is $3B$, giving $j=6$. In characteristic 13 it is $20A^3+15B^2$, giving $B^2=3A^3$ and $j=1728\cdot4/(4+81)=5$ in $\mathbb F_{13}$. For a nonsingular curve in the latter case $A\ne0$. The vanishing criterion is the usual Hasse--Deuring criterion \cite{silverman}.

If $p\nmid f$ and $p$ is nonsplit, all the $j$-values lie in a single relative disk of depth at least $1/2$ as above. For $k=4$, the strict thresholds are $1/4$ at 7 and $1/10$ at 13. Thus four conjugates suffice, for every $D$ in this case.

For a split prime $p\nmid f$, the alternative certificate is archimedean. For a primitive reduced form $[a,b,c]$ of discriminant $-D$ write $j_a=j((-b+i\sqrt D)/(2a))$ for its associated conjugate (the choice of $b$ is understood). The classical cusp bound in \cite[\S2]{blp} gives
\[
       \left||j_a|-\exp(\pi\sqrt D/a)\right|\le2079.
\]
For $a<b$ two leading coefficients, the tests
\[
 D\ge16a^2,\qquad 9D(b-a)^2\ge16a^2 b^2                                    \tag{1}
\]
imply $|j_a|>16|j_b|$. Indeed the two exponentials $U,V$ satisfy
$U\ge e^{4\pi}>34\cdot2079$ and $U/V\ge e^{4\pi/3}>32$, implying $U-2079>16(V+2079)$.
If four nonzero conjugates have successive moduli increasing by factors at least 16, their generalized Vandermonde is invertible. For increasing moduli $r_j$ and exponents $n_j$, scale successive columns with factors in ratio $(r_jr_{j+1})^{-(n_{j+1}-n_j)/2}$. Moving away from the diagonal in a row now decreases the entry modulus by a factor at least four per step. The other entries sum to at most $2\sum_{m\ge1}4^{-m}=2/3$ times the diagonal. This yields strict diagonal dominance. This argument works for any number of rows.

For $D\ge 4p^6$, the numbers $a=1,p,p^2,p^3$ are available as leading coefficients. At odd $p$ lift a nonzero root of $b^2\equiv -D\bmod p$ to the required power, choose the parity of $b$ for $4\mid b^2+D$, and take $|b|\le a$. Then $c=(b^2+D)/(4a)\ge a$ and the form is primitive. The principal form handles $a=1$. The consecutive pairs pass (1) at $4p^6$, hence thereafter, for both 7 and 13.

It remains possible that $7\cdot13\mid f$. For every $t=1,7,13,91$ construct a form with $a=t^2$. Put $b=tv$ with $v\equiv D\bmod2$, and choose $v$ so that $v^2+D/t^2\not\equiv0\bmod p$ for every $p\mid t$. At most two residues are forbidden at each odd prime. The Chinese remainder theorem permits $|v|\le t$. Then $c=(v^2+D/t^2)/4$ is integral, $\gcd(a,b,c)=1$, and if $D\ge4\cdot91^4$ then $c\ge a$. The successive entries of $1,49,169,8281$ pass (1) at this last cutoff. Signs on the boundary of the reduced fundamental domain may be normalized without changing $a$ or the modulus.

\subsection{The exhaustive check}
Set $A=470596=4\cdot7^6$, $B=19307236=4\cdot13^6$, and $C=274299844=4\cdot91^4$. For odd $p$, $p\mid f$ is equivalent to $p^2\mid D$. The preceding argument leaves only $D\equiv0,3\bmod4$ in three disjoint ranges:
\[
 3\le D<A,\qquad A\le D<B\ \text{with }49\mid D,\qquad
 B\le D<C\ \text{with }8281\mid D.
\]
A Python verifier checks their 442,904 entries. It first applies nonsplit reduction at 7 or 13, and then seeks four reduced leading coefficients up to 151 passing (1). For the surviving 51 it enumerates every reduced primitive form. It completes the following disjoint count.
\begin{center}
\begin{tabular}{lr}\toprule
Certificate or reason for discharge&Discriminants\\\midrule
Nonsplit prime and four conjugates whenever $h\ge4$&288,279\\
Four integer-tested leading coefficients&154,574\\
Class number below four&11\\
Complex determinant by interval arithmetic&28\\
Class-polynomial disk by integer arithmetic&12\\\midrule
Total&442,904\\\bottomrule
\end{tabular}
\end{center}
For the complex test order four disjoint moduli $0<r_0<\cdots<r_3$. Divide a permutation term of the determinant by the diagonal term in absolute value. If $g_i=n_{i+1}-n_i$, the quotient is
\[
 \prod_{i=0}^2\left(\frac{\prod_{j:\sigma(j)>i}r_j}{\prod_{j>i}r_j}\right)^{g_i}.
\]
Each bracket is at most one by the ordering. It suffices to check that the sum of the 23 other terms with exponents $0,1,2,3$ is strictly smaller in absolute-term sum than the diagonal magnitude. The code uses upper endpoints for the off-diagonal products and lower endpoints for the diagonal.

For the local test put $H(X)=\prod_j(X-j)$, the monic ring class polynomial of the discriminant in question (the product runs over its $j$-values).
Write $H(p^u(a+Y))=\sum_i c_iY^i$ with $p\nmid a$. The number of roots with
$v(Y)>r$ equals the \emph{smallest} index minimizing $v_p(c_i)+ir$, omitting zero coefficients. This follows by factoring over $\mathbb C_p$ and taking the initial polynomial at slope $-r$: a root of valuation $>r$ contributes a factor whose initial form is $Y$, one of valuation $<r$ a nonzero constant, and one on the boundary a linear factor with nonzero constant. Equivalently this is the Newton polygon theorem. Lemma~\ref{lem:disk} now applies with center $p^u a$.
For example, at $\Delta=-131$ the ascending coefficients of $H(8+Y)$ have 17-adic valuations $(4,3,2,1,0,0)$. At $r=1/14$ the least minimizing index is 4. This one row rules out every choice of four exponents, despite the presence of a fifth conjugate outside the disk.

The 12 local rows used by the independent proof are part of a larger reconciliation with the prior report. We checked every one of its 79 candidates directly. The disk certificates group as follows:
\[
\begin{array}{c|rrrrr}\text{prime}&2&3&7&13&17\\\hline
\text{discriminants}&11&39&22&5&2\end{array}
\]
In 28 of the 29 listed degree-five cases the same test at $k=5$ yields $s=6$. The remaining case of that list is $\Delta=-131$, where this proof gives $5\le s\le6$.
The data and checks are in \texttt{checks/disks.csv}, \texttt{local\_disks.py} and \texttt{global\_sieve.py} (the scripts are in \texttt{checks/}).

The class-polynomial verifier recomputes all 79 tables, encompassing 345 distinct conjugates counted across discriminants. For each reduced form with $q=\exp(2\pi i\tau)$, it uses
\[
 j(\tau)=\frac{\big(1+240\sum_{n\ge1}\sigma_3(n)q^n\big)^3}
                 {q\prod_{n\ge1}(1-q^n)^{24}}.
\]
All operations enclose real and imaginary parts using integers with common denominator $2^{1000}$. Reducedness gives $|q|<1/200$. After $N=170$ terms the Eisenstein error is bounded by $960(N+1)^3/200^{N+1}$, and the omitted product differs from 1 by at most $2/(199\cdot200^N)$. These bounds follow from $\sigma_3(n)\le2n^3$ and the geometric series. Multiplication of the linear factors encloses a unique integer for each coefficient, which identifies the exact class polynomial by CM integrality. The same bounds enclose the complex determinants. This completes the finite part of Theorem~\ref{th:four}.

\subsection{An exponential restriction at discriminant \texorpdfstring{$-131$}{-131}}
The class polynomial has degree five. Its constant coefficient is
$2^{81}17^3 23^3$, and reduction modulo $17$ gives $X(X-8)^4$.
Consequently one root $b$ has $v_{17}(b)=3$ and the other four,
$u_1,\ldots,u_4$, satisfy $v_{17}(u_i-8)\ge1/4$.
Indeed, the simple zero lifts by Hensel's lemma, and the valuation of
the product of all roots gives its valuation. For a root reducing to
$8$, if $0<v_{17}(u_i-8)<1/4$ the fourth-degree term in the expansion
about $8$ would be the unique term of minimum valuation.

Suppose that a polynomial of precisely five terms vanishes at these
roots. Removing its smallest monomial factor gives exponents
$0=e_0<e_1<e_2<e_3<e_4=D$. Then some $j\in\{1,2,3,4\}$ satisfies
\[
 3e_j\ \leq\!
 \sum_{\substack{1\leq k\leq4\\ k\ne j}}
    \big(v_{17}(e_k-e_j)-v_{17}(e_k)\big).
 \tag{Q}\label{eq:quinticgap}
\]
In particular $D\ge17^{e_1}$, and some difference $e_k-e_j$ with
$k\ge1$, $k\ne j$, is divisible by $17^{e_j}$.

To prove (Q), write $w_i=\log(u_i/8)$, so that
$v_{17}(w_i)\ge1/4$ and the $w_i$ are distinct. For any four distinct
nonnegative integers $n_l$ the determinant $\det(u_i^{n_l})$ has
valuation
\[
 B + \sum_{l<m}v_{17}(n_m-n_l),
 \qquad B=\sum_{i<h}v_{17}(w_h-w_i).
\]
Here is a direct valuation check. Remove the unit column factors
$8^{n_l}$ and expand $\det(\exp(w_i n_l))$ by Cauchy--Binet.
For indices $0\le t_0<\cdots<t_3$, dividing the two alternants by
their ordinary Vandermondes gives integral Schur polynomials of
degree $b'=\sum_{i=0}^3(t_i-i)$ in $w$ and in $n$, respectively.
Their combined valuation is at least $b'/4$. The factorials in the
denominator cost at most $b'/14$: for $t_i<17$ the valuation is zero;
for $t_i\ge17$, Legendre's bound gives
$v_{17}(t_i!)\le(t_i-1)/16\le(t_i-i)/14$.
Thus the leading term $(0,1,2,3)$ is a unit after division, and all
later terms have positive valuation tending to infinity. This proves
the asserted exact valuation.

The five-by-five alternant on the roots is singular because its
columns satisfy the rational relation. Expand it along the row for
$b$. The term indexed by $j$ has valuation $3e_j+B$ plus the
valuation of the Vandermonde on the other four exponents. Some term
with $j>0$ must have valuation at most that of term $0$.
Cancel the pairs common to the two minors to obtain (Q). Each of its
three summands is at most $\log_{17}D$.

\paragraph{A certified degree cutoff.}
For $D\le1000$ the anchor exponent $e_j$ in (Q) must equal $1$ or
$2$, since $17^3>1000$. The program
\texttt{checks/degree\_131.py} enumerates each triple of remaining
positive exponents satisfying (Q), sorted into buckets by
$v_{17}(e_k-e_j)$. It finds $179{,}365$ admissible triples for anchor
$1$ and one for anchor $2$. For each, reduce the five powers of $X$
modulo the monic class polynomial and then modulo $1{,}000{,}000{,}007$.
The powers $1$ and $X^{e_j}$ are coordinate vectors, so a
three-by-three determinant on the remaining coordinates tests full
rank. Every tested determinant is nonzero modulo that integer.
It follows that every five-term relation would have normalized degree
at least $1001$. The class polynomial provides a six-term relation,
and the existence of a five-term one of higher degree remains open
in this analysis.

\section{How the minimum support grows}

Let $D=|\Delta|$, $h=h(\Delta)$, and let $s(\Delta)$ denote the minimum
number of terms of a nonzero rational polynomial vanishing at a singular
modulus of discriminant $\Delta$. Write $\Delta=f^2 d$, where $d$ is a
fundamental discriminant. We obtain two uniform bounds as $D\to\infty$:
\[
 s(\Delta)\gg_{\epsilon}D^{1/4-\epsilon}\quad(\epsilon>0),
 \qquad
 s(\Delta)\gg\sqrt h\gg\frac{D^{1/4}}{\sqrt{\log\log(3D)}}
 \quad\text{using Dirichlet GRH}.
\]
The first bound uses Siegel's ineffective constant. The second improves the
$D^{1/4}/\sqrt{\log D}$ bound in the preceding ProveIt manuscript. Both
include nonmaximal orders.

\paragraph{Class numbers and reduction.}
The ring class number formula (Cox, \emph{Primes of the Form
$x^2+ny^2$}, Thm.~7.24) gives
\[
 h=\frac{h(d)f}{[\mathcal O_K^\times:\mathcal O_f^\times]}
       \prod_{\ell\mid f}\left(1-\frac{\chi_d(\ell)}\ell\right)
   \ge \tfrac13 h(d)\varphi(f),
\]
where $K=\mathbb Q(\sqrt d)$ and $\chi_d=(d/\cdot)$.
The totient bound $\varphi(f)\gg f/\log\log(3f)$
(Rosser--Schoenfeld, \emph{Illinois J. Math.} 6 (1962), Thm.~15)
and $h(d)\gg_\epsilon |d|^{1/2-\epsilon}$ give
$h\gg_\epsilon D^{1/2-\epsilon}$. For the field
bound see \href{https://doi.org/10.4064/aa-1-1-83-86}{Siegel (1935)}.

Under Dirichlet GRH the conductor factors can be retained in the analytic
class number formula. Put $y=(\log(3D))^6$. The usual GRH estimate
$\sum_{p\le t}\chi_d(p)\log p\ll\sqrt t\log^2(3Dt)$, followed by
partial summation, gives
\[
 \left|\sum_{p>y}\frac{\chi_d(p)}p\right|
 \ll \frac{\log^2(3Dy)}{\sqrt y\log y}
       +\int_y^\infty\frac{\log^2(3Dt)}{t^{3/2}\log t}\,dt
 \ll 1 .
\]
For the last inequality split $\log^2(3Dt)\le
2\log^2(3D)+2\log^2 t$ and use this choice of $y$.
In the logarithm of the Euler product the terms of prime-power exponent at
least two have bounded absolute sum. Also
$\sum_{p\mid f,\,p>y}1/p\le\log(2f)/(y\log y)$. Hence
\[
 \log\!\left(L(1,\chi_d)\prod_{p\mid f}
       (1-\chi_d(p)/p)\right)
 =\sum_{\substack{p\le y\\p\nmid f}}\frac{\chi_d(p)}p+O(1)
 \ge-\log\log y-O(1)
\]
by Mertens' prime reciprocal sum. The Euler expansion at $1$ follows by
taking the limit from the right, justified by the displayed tail estimate.
The class number formula now gives $h\gg\sqrt D/\log\log(3D)$ for all
orders. This combines the finite Euler factors before estimating the
$L$-value; for explicit bounds on the primitive value see
\href{https://arxiv.org/abs/1309.3595}{Lamzouri--Li--Soundararajan,
Thm.~1.5}.

At an inert $p\nmid f$ all $h$ conjugates reduce to supersingular
invariants, and the ring class field is unramified at $p$.
There are $O(p)$ supersingular invariants. For example, every elliptic curve
in characteristic $p>3$ has a Legendre model over the algebraic closure, and
its supersingular parameter is a root of the nonzero Hasse polynomial
$\sum_{i=0}^{(p-1)/2}\binom{(p-1)/2}{i}^{\!2}\lambda^i$.
For $p\equiv1\pmod3$, the invariant $0$ is ordinary: the coefficient of
$x^{p-1}$ in $(x^3+1)^{(p-1)/2}$ is
$\binom{(p-1)/2}{(p-1)/3}\not\equiv0\pmod p$.
If $M$ conjugates have the same nonzero reduction, divide them all by one
of these conjugates. The scaled points $z$ satisfy $v_p(z-1)\ge1$.
For $m=\min(M,\lfloor p/2\rfloor)$ with $m\ge2$ and $p\ge7$, the
local determinant estimate applies to any $m$ of them, since
$1/(p-m+1)<1$. Thus $s(\Delta)>m$. Common scaling multiplies each
power column by a nonzero scalar.

\paragraph{A bound independent of zero-location theorems.}
Set $X=D^{1/4}$. Among the primes $X/2<p\le X$, those congruent to $1$
modulo $12$ and coprime to $D$ number $\gg X/\log X$. This is the
unconditional prime number theorem in the fixed progression, with at most
$\log D/\log 2$ primes removed; see
\href{https://arxiv.org/abs/1802.00085}{Bennett--Martin--O'Bryant--Rechnitzer}.
If one is inert, at least $h/O(X)$ conjugates have one common nonzero
reduction. The disk bound gives $s(\Delta)\gg\min(X,h/X)$, which
suffices by the class number estimate above.

Otherwise all the indicated primes split. Each gives a primitive reduced
form of leading coefficient $p$: choose $|b|\le p$ with
$b^2\equiv \Delta\pmod{4p}$, and take $c=(b^2+D)/(4p)>p$.
The corresponding conjugate satisfies
$\log |j|=\pi\sqrt D/p+o(1)$ uniformly by the cusp bound
$\big||j|-\exp(\pi\sqrt D/p)\big|\le2079$.
Two different such primes give exponents separated by at least
$\pi\sqrt D/X^2=\pi$. Hence consecutive moduli, in increasing order,
have ratio $>16$ for large $D$. To see that $t$ such conjugates exclude
every $t$-term relation, order the moduli $r_i$ and exponents $e_j$
increasingly. Scale column $j$ of the power matrix by $c_j>0$ with
$c_{j+1}/c_j=(r_jr_{j+1})^{-(e_{j+1}-e_j)/2}$. Moving away from the
diagonal in each row decreases absolute entries by factors of at least
four. The off-diagonal sum is at most $2/3$ of the diagonal entry.
Strict diagonal dominance gives independence. Thus
$s(\Delta)\gg X/\log X$ in this case.

\paragraph{A sharper bound using the proved Dirichlet GRH.}
Take $Y=\sqrt h$. We seek an inert prime $Y<p\le 2Y$, coprime to $f$,
where at least half the conjugates are units. Dirichlet GRH gives
\[
 \sum_{\substack{Y<p\le 2Y\\\chi_d(p)=-1,\ p\nmid f}}\log p
  = Y/2+O\big(\sqrt Y\log^2(3DY)+\log(3D)\big)\gg Y.
\]
Indeed, subtract the weighted sums for $\chi_d$ from the ordinary prime
sums and divide by two. The ramified and conductor primes cost $O(\log D)$.
The standard GRH error also follows by subtracting the quadratic prime
ideal estimate of
\href{https://arxiv.org/abs/1312.4463}{Greni\'e--Molteni, Cor.~1.2}
from twice the rational prime estimate and halving, with
$O(\sqrt{Y}\log^2(3Y)+\log D)$ for prime powers and ramification.
Here $Y\gg D^{1/4}/\sqrt{\log\log(3D)}$, so the error is $o(Y)$.
In particular this selection works for $d=-3$ as well.

Write $N=|H_\Delta(0)|$, the absolute product of the conjugates.
For large $D$, $N$ is a nonzero integer. We claim that for every $\eta>0$,
\[
                         \log N\ll_\eta D^{1/2+\eta}.             \tag{*}
\]
For a reduced form $[a,b,c]$, we have $a\le\sqrt{D/3}$ and
$\log^+ |j|\ll1+\sqrt D/a\ll\sqrt D/a$. The number of possible
$b$ with $|b|\le a$ is $\ll 2^{\omega(a)}\sqrt{\gcd(a,D)}$.
Indeed, modulo an odd prime power $\ell^k$, a square congruence to $\Delta$
has at most $2\sqrt{\gcd(\ell^k,D)}$ solutions: extract the largest common
even valuation and count unit square roots, or count multiples of
$\ell^{\lceil k/2\rceil}$ when $\ell^k\mid D$. At $2$ there are at most
four unit square roots. Replacing $a$ by $4a$ costs an absolute factor.
Apply the Chinese remainder theorem and sum the resulting upper bound:
\[
 \sum_{a\le\sqrt D}\frac{2^{\omega(a)}\sqrt{\gcd(a,D)}}a
 \le \sum_{r\mid D}\frac{2^{\omega(r)}}{\sqrt r}
            \sum_{n\le\sqrt D}\frac{2^{\omega(n)}}n
 \ll \tau(D)^2\log^2(3D).
\]
Here $\tau$ counts divisors, $2^{\omega(n)}\le\tau(n)$, and
$\tau(D)\ll_\eta D^\eta$ for every $\eta>0$. This proves $(*)$.

At the chosen inert primes the normalized valuations of the integral
conjugates are nonnegative integers. Each conjugate reducing to zero
contributes at least one to $v_p(N)$. If more than half reduced to
zero at each such prime, we would have $\log N\gg hY=h^{3/2}$,
contrary to $(*)$ and the lower bound for $h$ (take, for example,
$\eta=1/16$). At some prime, therefore, at least $h/2$ conjugates
are units. Among its $O(p)$ supersingular invariants, one nonzero
invariant receives $\gg h/p\gg\sqrt h$ of the conjugates. Since
$p\asymp\sqrt h$, the disk estimate gives $s(\Delta)\gg\sqrt h$,
as asserted.

\section{A quantitative answer to Suzuki's density question}
\label{sec:suzuki-density}
For $x>1$ let $E(x)=\sum_{n\le x}'\Lambda(n)/\sqrt n-2\sqrt x$.
The prime on the sum gives half weight to its last term when $x$ is a prime power.
Suzuki explicitly asks for the logarithmic density of $E(x)\le0$ in
\S1.1 after (6) of \emph{On variants of Chebyshev's conjecture}
(\href{https://arxiv.org/html/2411.07436v3}{corrected text}).
His Theorem~1 proves a sign statement for an integral of $E$.
The incoming report \texttt{grh-transfers.zip}, file \texttt{races.tex},
compares primes in distinct residue classes with an additional logarithmic smoothing.
Here we bound how often the \emph{unsmoothed} error itself is negative.

\paragraph{Theorem.}
Put $p(a,L)=L^{-1}|\{t\in[a,a+L]:E(e^t)>0\}|$, where $|\cdot|$
denotes Lebesgue measure. Using RH,
\[
 0<\liminf_{L\to\infty}\inf_{a\ge1}p(a,L)
 \ \le\ \limsup_{L\to\infty}\sup_{a\ge1}p(a,L)<0.00323 .
\]
In particular the lower logarithmic density of $E(x)<0$ exceeds
$99.677\%$, and its upper logarithmic density is strictly below one.
The placement of the logarithmic interval is arbitrary. Global simplicity
and linear independence of zero ordinates are both unnecessary.
The ordinary logarithmic density at this particular threshold is still undetermined by this argument.

\paragraph{The frequencies and their mean square.}
Write $\gamma$ for a distinct positive ordinate of a nontrivial zeta zero and
$m_\gamma$ for its multiplicity. Suzuki's equations (36) and (47) give, under RH,
\begin{align*}
 E(e^t)&=-c+F(t)+r(t),&
 c&=\tfrac12(\pi/2+\log(8\pi)+\gamma_{\rm E}),\\
 F(t)&=-2\lim_{T\to\infty}\sum_{0<\gamma\le T}
                 m_\gamma\frac{\sin(\gamma t)}{\gamma},&
 r(t)&=\sum_{j\ge1}\frac{e^{-(2j+1/2)t}}{2j+1/2}.
\end{align*}
Here $\gamma_{\rm E}$ is Euler's constant and $r(t)$ tends exponentially to zero.
This uses the same weighted-prime-sum/zero-residue correspondence as the
repository's \texttt{New/Reduction/Zeros/PrimeError}, with Suzuki's
weight in place of its one-sided remainder. The geometric determinant
inequality enters this particular deduction through its final RH theorem.

We give the convergence argument because pointwise explicit formulas alone
do not justify the density estimate. Define the Stepanov norm
$\|f\|_{S^2}=\sup_b(\int_b^{b+1}|f(t)|^2dt)^{1/2}$ on the real line.
For a finite exponential sum $P(t)=\sum A_\lambda e^{i\lambda t}$, put
$B_k=\sum_{k\le\lambda<k+1}|A_\lambda|$. Then
\[
                \|P\|_{S^2}^2\ \ll\ \sum_{k\in\mathbb Z} B_k^2 .
\]
Indeed, $W(u)=(\sin(u/2)/(u/2))^2$ dominates a positive constant on
$[-1/2,1/2]$ and its Fourier transform with angular frequencies is bounded
and supported on $[-1,1]$. Integrate $|P|^2$ against each translate of $W$.
Only equal or adjacent unit frequency blocks contribute, and
$2B_kB_{k+1}\le B_k^2+B_{k+1}^2$.
The Riemann--von Mangoldt estimate gives $O(\log(2+k))$ zeros with
multiplicity per unit block at positive height $k$.
For tails of the displayed sine series this gives
$B_k\ll\log(2+|k|)/(1+|k|)$ away from the origin.
The squared bound is summable. The truncations therefore have an $S^2$
limit on the real line, agreeing almost everywhere for $t>0$ with the
explicit formula by local $L^2$ convergence and pointwise convergence.
This is a direct instance of the Stepanov approach of
Kaczorowski and Ramar\'e, \emph{Acta Arith.} 106 (2003), 277--297
(\href{https://doi.org/10.4064/aa106-3-6}{paper}).

For a finite truncation $P$, the means of $P^2$ and $(P_+)^2$ over
intervals of length $L$ converge uniformly in their starting point.
In fact the phase flow is uniformly distributed in its compact orbit closure,
with Haar measure there. This follows on characters by integration of
$e^{i\lambda t}$, and on continuous functions by uniform approximation.
Conjugation on the phase group preserves Haar measure and sends the sine
sum to its negative. Thus the second mean is half the first, including
when ordinates satisfy rational relations.
The positive-part map is $1$-Lipschitz. On any interval of length $L\ge1$,
the normalized $L^2$ norm of the difference between $P$ and $F$ is at most
$\sqrt{1+1/L}\,\|P-F\|_{S^2}$. Taking the truncation limit and using
orthogonality of distinct frequencies yields the uniform mean
\[
          \lim_{L\to\infty}\frac1L\int_a^{a+L} (F(t)_+)^2dt
          =\frac V2,\qquad V=2\sum_{\gamma>0}\frac{m_\gamma^2}{\gamma^2}.
\]
For any $0<\eta<c$, the inequality $E(e^t)>0$ forces
$F(t)>c-\eta$ after a fixed starting time. The initial segment contributes
$O_\eta(1/L)$ uniformly for $a\ge1$.
Markov's inequality, followed by $\eta\downarrow0$, gives $V/(2c^2)$
as an upper bound for the upper density.

\paragraph{An explicit multiplicity allowance.}
Put $H=3\cdot10^{12}$ and $S=2+\gamma_{\rm E}-\log(4\pi)$.
The Hadamard product for $\xi$, on subtracting its logarithmic derivatives at
$1$ and $0$, gives
\[
 S=2\xi'(1)/\xi(1)=2\sum_{\gamma>0}\frac{m_\gamma}{1/4+\gamma^2}.
\]
Here the first evaluation follows from the Laurent expansion of $\zeta$ at
one and the gamma factor. For $u=(1/4+\gamma^2)^{-1}$ one has $u\le S/2$
and $\gamma^{-2}=u/(1-u/4)$. Thus
$2\sum m_\gamma/\gamma^2\le S/(1-S/8)$.
Platt and Trudgian verified simplicity through $H$
(\href{https://doi.org/10.1112/blms.12460}{Bull. London Math. Soc. 53 (2021)}).
At larger heights the zero-count estimate of Hasanalizade, Shen and Wong
(\href{https://arxiv.org/abs/2107.06506}{Corollary 1.2}) gives
\[
 \left|N(t)-\frac{t}{2\pi}\log\frac{t}{2\pi e}\right|
 \le 0.1038\log t+0.2573\log\log t+9.3675 <\log t \quad(t\ge H).
\]
Consequently $N(t)\le t\log t$ and its jump at $\gamma>H$ is at most
$2\log\gamma$. The extra multiplicities in $V$ contribute at most
\[
 4\int_{(H,\infty)}\frac{\log t}{t^2}\,dN(t)
 \le 8\int_H^\infty \frac{\log^2t}{t^2}\,dt
 =\frac{8}{H}\big((\log H)^2+2\log H+2\big).
\]
Integration by parts drops a nonpositive lower boundary term.
It follows that $V<0.04647$, $c>2.686$, and $V/(2c^2)<0.00323$.

\paragraph{Why the other sign has positive lower density.}
The classical Littlewood oscillation theorem makes
$e^{-t/2}(\psi(e^t)-e^t)$ unbounded above and below. Its zero series
has coefficients $-1/(1/2+i\gamma)$, differing from those of $F$ by
$O(\gamma^{-2})$ per zero. The difference series is uniformly bounded,
as are the remaining terms at $t\ge1$. Thus $E(e^t)$ also takes arbitrarily
large positive and negative values (as Suzuki observes in \S1.1).
The difference between the right-continuous and the half-weight convention
tends to zero. Between consecutive prime powers $E(x)$ is smooth, and at
a jump its one-sided values bracket the half-weight value.
Choose a sufficiently late interval $I$ of length at most one and
$\eta>0$ such that $F>c+2\eta$ almost everywhere on $I$.

The $S^2$ limit of finite trigonometric polynomials has relatively dense
$S^2$ almost periods: approximate it by one polynomial, then use compactness
of that polynomial's phase group to obtain relatively dense uniform almost
periods. For every such $\tau$ with
$\|F(\cdot+\tau)-F\|_{S^2}<\eta\sqrt{|I|/2}$, at least half of $I+\tau$
has $F>c+\eta$, by Chebyshev's inequality. These translates occur with
bounded gaps. More explicitly, if every shift interval of length $D$ contains
such a $\tau$, each consecutive block of length $D+|I|$ contains one
complete translated interval. Since $r(t)>0$, this proves the stated
strictly positive uniform lower density. Finally $E(x)=0$ on a countable
set because it is strictly decreasing between prime powers. This permits
either sign convention at the density threshold.

\section{An exact change of sign in Suzuki's average over moduli}
\label{sec:fixed-suzuki}
Suzuki's equation (29), also printed in the \href{https://doi.org/10.1007/s11139-025-01289-y}{journal correction}, concerns
\[
 A_Q(x)=\sum_{q=3}^Q\left(
 \varphi(q)\sum_{\substack{n\le x\\n\equiv1\bmod q}}
       \frac{\Lambda(n)}{\sqrt n}\log\frac{x}{n}-4\sqrt x\right).
\]
His Theorem~7 gives negativity when $Q\ge x$. The paragraph after (30)
asks about a large fixed $Q$. The order of limits changes the answer.
Use the repository's Dirichlet GRH theorem throughout this section.

\paragraph{Theorem.}
For each fixed integer $3\le Q\le138$, $A_Q(x)$ is eventually negative.
For each fixed $139\le Q\le2{,}000{,}000$, it is eventually positive.
More generally, for any fixed $Q\ge139$ it is eventually negative if and
only if some primitive Dirichlet $L$-function of conductor at most $Q$
vanishes at $1/2$. In the other case it is eventually positive.
The starting value of $x$ may depend on $Q$.

\paragraph{The central residue.}
Set $Z_q(s)=\prod_{\chi\bmod q}L(s,\chi)$, $m_q=\operatorname{ord}_{1/2}Z_q$,
and define $C(q)$ by $Z_q'/Z_q(1/2+z)=m_q/z+C(q)+O_q(z)$.
Put $M(Q)=\sum_{q=3}^Q m_q$ and $B(Q)=\sum_{q=3}^Q C(q)$.
Character orthogonality and the twice-integrated explicit formula give
\[
 A_Q(x)=-\tfrac12 M(Q)(\log x)^2-B(Q)\log x+O_Q(1).
\]
In detail, the Mellin integrand for the $q$th inner weighted sum, including
its factor $\varphi(q)$, is $-(Z_q'/Z_q)(1/2+z)x^z/z^2$.
The pole at $z=1/2$ contributes $4\sqrt x$. The residue at zero yields
the two displayed powers and a constant. The noncentral critical zeros
yield a bounded oscillating sum because their denominators are
$|\Im\rho|^2$ and the modulus is fixed.
For primitive factors this is precisely the integrated explicit formula
used in Suzuki, \S\S3.1, 4.1 and 5.1
(\href{https://arxiv.org/html/2411.07436v3}{paper}).
For each Euler factor deleted on induction, expand its logarithmic
derivative in powers of $p^{-1/2-z}$. The coefficients at $z=0$ are
geometrically decreasing, so termwise summation with $\log(x/p^j)$
adds its stated linear term and a bounded remainder. This also justifies
the formula for the principal character.

\paragraph{An elementary coefficient valid at a central zero.}
Let $h=\varphi(q)$, $K=\mathbb Q(\zeta_q)$ and $D=|\operatorname{disc}K|$.
For $p^a\Vert q$ put $d=q/p^a$, $e=\varphi(p^a)$,
$f=\operatorname{ord}_d(p)$ with $f=1$ when $d=1$, and $g=h/(ef)$.
Euler products give $Z_q=\zeta_K\prod_{p\mid q}(1-p^{-fs})^g$.
The completed Dedekind zeta function, up to a fixed positive multiple,
is $D^{s/2}(2\pi)^{-hs/2}\Gamma(s)^{h/2}\zeta_K(s)$
and is symmetric about $1/2$. Its logarithmic derivative therefore has
zero constant Laurent coefficient at $1/2$, including at a zero.
The standard cyclotomic discriminant formula
$\log D=h(\log q-\sum_{p\mid q}\log p/(p-1))$
and $\Gamma'(1/2)/\Gamma(1/2)=-\gamma_{\rm E}-2\log2$ now give
\[
 C(q)=\frac{h}{2}\left(\log(8\pi)+\gamma_{\rm E}-\log q+
                         \sum_{p\mid q}\frac{\log p}{p-1}\right)
      +\sum_{p^a\Vert q}\frac{\varphi(q/p^a)\log p}{p^{f/2}-1}.
\]
Thus both the cutoff computation and the large-$q$ sign estimate apply
irrespective of uncomputed central values.

\paragraph{Certification of the crossover.}
The verifier \path{checks/suzuki_fixed_q.py}
encloses $B(Q)$ for every $3\le Q<10^4$. It finds
$45.431<B(138)<45.490$ and $-29.828<B(139)<-29.767$,
with $B(Q)>0$ throughout $3\le Q\le138$ and $B(Q)<0$ throughout
$139\le Q<10^4$. It uses rational remainder bounds for the logarithm
series after binary range reduction and integer square roots for the
ramified Euler factors. Machin's identity, the harmonic sum through 1024 with a trapezoid-rule
remainder, and the same logarithm series certify
$3.80138<\log(8\pi)+\gamma_{\rm E}<3.80140$.

For $q\ge10^4$, each $C(q)$ is negative. Indeed
\[
 \sum_{p\mid q}\frac{\log p}{p-1}
 \le\tfrac14\log q+\tfrac34\log2+\tfrac14\log3
 <\tfrac14\log q+0.8.
\]
Only the primes $2,3$ require the two corrections.
As $p^f\ge d+1$, we have
$p^{f/2}-1\ge(\sqrt2-1)\sqrt d$, also for $d=1$.
Since $e\ge p^a/2$ and $p^a\ge2$, the last sum in $C(q)$ divided by
$h$ is at most $(2+\sqrt2)\log q/\sqrt q$.
Consequently $C(q)/h<2.31-(3/8-0.035)\log q<0$ in this range.
It follows that $B(Q)<0$ for every $Q\ge139$.

The central orders are nonnegative. Imprimitive Euler factors do not
vanish at $1/2$, so $M(Q)>0$ exactly when a primitive character of
conductor at most $Q$ gives a central zero (the zeta value at $1/2$ is
nonzero). Platt's independent computation (Theorem 9.2 in the author manuscript)
verifies central nonvanishing
through modulus $2{,}000{,}000$
(\href{https://doi.org/10.1090/mcom/3077}{Math. Comp. 85 (2016), 3009--3027}).
Its central-value bound is distinct from the finite-height GRH bound in
the same paper. Substitution in the fixed-$Q$ expansion proves the theorem.

\paragraph{The second cutoff outside a set of finite logarithmic measure.}
Suzuki's equation (30) uses the sum
\[
 U_Q(x)=\sum_{q=3}^Q\varphi(q)
 \sum_{\substack{n\le xe^2\\n\equiv1\bmod q}}
 \frac{\Lambda(n)}{\sqrt n}\log\frac{x}{n}.
\]
For every fixed $Q$ its sign agrees with the eventual sign of $A_Q$ outside
a set $\mathcal E_Q\subset[2,\infty)$ with
$\int_{\mathcal E_Q}dx/x<\infty$.
To see this, put $y=xe^2$ and
$D_Q(y)=\sum_{q=3}^Q\varphi(q)\sum_{n\le y,\,n\equiv1\bmod q}
\Lambda(n)/\sqrt n$. The unsmoothed explicit formulas (36), (43), (44)
give
\[
 D_Q(e^t)=2(Q-2)e^{t/2}-M(Q)t+R_Q(t),
 \qquad\sup_{b\ge1}\int_b^{b+1}|R_Q(t)|^2\,dt<\infty.
\]
Indeed each noncentral zero contributes a frequency with coefficient of
absolute value $1/|\Im\rho|$ per unit of multiplicity. The fixed-conductor
zero count in a unit block is $O_Q(\log(2+|k|))$. The Fourier cutoff estimate
for unit intervals in Section~\ref{sec:suzuki-density} thus applies to the
finite family. The finitely many deleted Euler factors contribute bounded
terms through geometrically decreasing prime-power sums. Endpoints at
prime powers form a countable exceptional set.
From $U_Q(x)=A_Q(y)+4(Q-2)\sqrt y-2D_Q(y)$ and $\log y=\log x+2$
we obtain, writing $u=\log x$,
\[
 U_Q(e^u)=-\tfrac12M(Q)u^2-B(Q)u+H_Q(u),
 \qquad \sup_{b\ge1}\int_b^{b+1}|H_Q(u)|^2\,du<\infty .
\]
Since $B(Q)\ne0$, the absolute value of the polynomial part exceeds
$c_Q k$ on $[k,k+1]$ for all sufficiently large integers $k$, for some $c_Q>0$.
The set of sign failures in this interval has measure $O_Q(k^{-2})$.
Summing proves the finite logarithmic measure assertion.

\paragraph{A growing range.}
There is also a direct answer for part of Suzuki's growing-$Q$ question.
The prime weights in (29) are nonnegative and
$\sum_{q\mid n-1}\varphi(q)=n-1$ for $n\ge2$. Thus, uniformly in integers $Q\ge3$,
\[
 A_Q(x)\le \sum_{n\le x}\Lambda(n)\sqrt n\log(x/n)-4(Q-2)\sqrt x
 = (4/9+o(1))x^{3/2}-4(Q-2)\sqrt x.
\]
Partial summation with the prime number theorem gives the last equality,
using $\int_0^1 u^{1/2}\log(1/u)\,du=4/9$.
For every fixed $\varepsilon>0$ this proves eventual negativity throughout
$Q\ge(1/9+\varepsilon)x$, directly from the positive-sum argument of
Suzuki's proof. The result holds independently of RH.

\paragraph{What remains.}
Eventual negativity for arbitrarily large \emph{fixed} $Q$ in (29)
is equivalent to the existence of a single central Dirichlet zero.
This identifies the precise obstruction for that question and gives the
opposite sign over the verified conductor range starting at $139$.
The transition in intermediate growing ranges, and eventual pointwise sign for
(30), require additional arguments.

\section{Balancing all prime-power races in a single multiplicative interval}
\label{sec:balancing}
Fix a positive integer $q$, put $G=(\mathbb Z/q\mathbb Z)^\times$ and
$d=|G|$, and let $\chi$ run over the characters of $G$. Write $\chi^*$
for its primitive character and $m_\chi$ for the order of vanishing of
$L(s,\chi^*)$ at $1/2$. For $a\in G$ define
\begin{align}
 \psi(x;q,a)&=\sum_{n\leq x,\ n\equiv a\pmod q}\Lambda(n),&
 b_a&=\frac{2}{d}\sum_\chi\overline{\chi(a)}m_\chi,\nonumber\\
 D_a(x)&=\frac{\psi(x;q,a)-x/d}{\sqrt x}+b_a,&
 D(x)&=(D_a(x))_{a\in G}.\label{eq:racevec}
\end{align}
The constants $b_a$ are real by conjugation. They incorporate possible
central zeros, so even characters with unknown central nonvanishing
are included. Using the von Mangoldt weight counts prime powers as well as primes.

\paragraph{Fixed-interval balancing theorem.}
Assume the upstream Dirichlet GRH. There exist $C_q>1$ and $r_q>0$
such that, for every sufficiently large real $X$,
\[
 \{y\in\mathbb R^d:\|y\|_2\leq r_q\}
 \ \subseteq\ \operatorname{conv}\{D(n): n\in\mathbb Z,\ X\leq n\leq C_qX\}.
 \tag{B}\label{eq:ballrace}
\]
Consequently every such $y$, including zero, is exactly a convex
combination of at most $d+1$ integer samples from that interval, with
the same sample points and weights for every residue class. Every
unit linear combination of the $d$ coordinates attains values at least
$r_q$ and at most $-r_q$ in the interval. The theorem gives existence
with constants depending on the fixed modulus. It requires neither
simplicity of zeros nor rational independence of their ordinates.

\subsection{A local averaging argument from the Lean proof}

The compact proofs of \texttt{IsPolar.eventually\_moment\_le} and
\texttt{IsPolar.exists\_window\_sub\_mul\_le} use high Laplace moments to
find a one-sided small value on $[t,t^2]$ in log coordinates.
They are in
\href{https://github.com/lsernau/grh-lean/blob/72fead102f663d5c7a407964a13dc0cf9d977207/grh-compact/src/New/Endgame/OneSidedCriterion.lean}{\texttt{New/Endgame/OneSidedCriterion.lean}}.
The following vector version explains how a common choice of positive
weights replaces a separate choice of a point for each direction.

\paragraph{Laplace averaging lemma.}
Let $E:[0,\infty)\to\mathbb R^d$ be measurable, with
$\|E(v)\|\leq M e^{Av}$, $A\geq0$. Suppose its coordinate Laplace
transforms $F(z)=\int_0^\infty E(v)e^{-zv}\,dv$, initially defined
for $\Re z>A$, extend holomorphically to $\Re z>0$ and a neighborhood of
$0$. For every $h>0$ there is $\delta>0$ such that, for all large $T$,
a probability density supported on $[T,(1+h)T]$ has $E$-average of
norm $O(e^{-\delta T})$. The same estimate holds for a convex
combination of at most $d+1$ values in the interval.

\paragraph{Proof.}
Take $c=1+h/2$. Choose $\lambda>A$ sufficiently large as specified
below, set $k=\lfloor\lambda cT\rfloor$, and let $V$ have density
$\lambda^{k+1}v^ke^{-\lambda v}/k!$ on $v>0$. The disk
$|z-\lambda|\leq\lambda$ meets the imaginary axis only at $0$.
Compactness gives $R>\lambda$ with $F$ holomorphic on a neighborhood
of the closed concentric disk of radius $R$. Cauchy's estimate yields
\[
 \|\mathbb E E(V)\|
 =\lambda^{k+1}\|F^{(k)}(\lambda)\|/k!
 \ll\lambda(\lambda/R)^k.
\]
For the tails, exponential tilting gives
\[
 \mathbb E[e^{AV};\ V\notin I]
 =\left(\frac{\lambda}{\lambda-A}\right)^{k+1}
      \Pr(W\notin I),\qquad I=[T,(1+h)T],
\]
where $W$ has the same shape $k+1$ and rate $\lambda-A$. Its mean
divided by $T$ tends to $c\lambda/(\lambda-A)$.
For all sufficiently large $\lambda$ this limit stays a fixed
positive distance from both endpoints $1,1+h$. The Chernoff bound for
a gamma variable at $u$ times its mean is
$\exp[-(k+1)(u-1-\log u)]$ for the corresponding tail. This follows
by optimizing its moment generating function.
The two endpoint ratios stay bounded away from $1$, so the probability
is at most $2e^{-\kappa_h(k+1)}$ with $\kappa_h>0$ independent of all
sufficiently large $\lambda,T$. Choose $\lambda$ large enough that
$\log(\lambda/(\lambda-A))<\kappa_h/2$. Truncating the density to $I$
thus changes its vector integral by $O(e^{-\delta T})$, and its mass
is $1-O(e^{-\delta T})$, after decreasing $\delta$ if necessary.
The normalized integral belongs to the closed convex hull. Approximate
it to $e^{-\delta T}$ by a finite convex combination and apply
Carath\'eodory's theorem.\hfill$\square$

For the character coordinate of \eqref{eq:racevec}, write
$\psi(x,\chi)=\sum_{n\le x}\chi(n)\Lambda(n)$,
$e_\chi=1$ for the principal character and $0$ otherwise, and
$s=z+1/2$. The Laplace transform of
$e^{-v/2}(\psi(e^v,\chi)-e_\chi e^v)+2m_\chi$ is
\[
 -\frac{1}{s}\frac{L'}{L}(s,\chi)-\frac{e_\chi}{s-1}
                +\frac{2m_\chi}{z}.
\]
The pole at $s=1$ cancels for the principal character. The residue
at $z=0$ before adding the final term is $-2m_\chi$. Finite Euler
factors removed from primitive $L$-functions have zeros only on
$\Re s=0$. Dirichlet GRH therefore gives the continuation required
by the lemma. Fourier inversion gives all the coordinates in
\eqref{eq:racevec}. The lemma already gives exponentially accurate
balance in each log-window $[T,(1+h)T]$.
For exact balance in a log-window of fixed length we use the
individual nonzero frequencies.

\subsection{Simultaneous distinction of the zeros}

\paragraph{Multiplicity rank theorem.}
For each fixed $q$ there is $K_q$ such that if $T$ is sufficiently
large and $K_q\log T\leq H\leq T$, the vectors
\[
 M(\rho)=\big(\operatorname{ord}_\rho L(s,\chi^*)\big)_{\chi\bmod q}
 \quad (0<\Re\rho<1,\ T<\Im\rho<T+H)
\]
span $\mathbb R^d$. Distinct zeros index the vectors; a zero absent
from one $L$-function contributes a zero entry. This assertion about
complex zeros holds unconditionally. Under Dirichlet GRH the same
vectors are indexed by distinct positive ordinates.

\paragraph{Proof.}
Choose a prime $p_a>2$ in each class $a\in G$ by Dirichlet's theorem.
Banks's short-interval Landau formula gives
\[
 A_{a\chi}:=\sum_{T<\Im\rho<T+H}p_a^\rho M(\rho)_\chi
  =-\frac{H\log p_a}{2\pi}\chi(a)+O_q(\log T).
\]
Here the sum is over the union of distinct nontrivial zeros. At these
primes $\chi^*(p_a)=\chi(a)$. The implied constant can be fixed for
the finite list of characters and chosen primes. The row-scaled
character table multiplying $H$ is invertible. Choose $K_q$ so its
smallest singular value times $H$ exceeds the operator norm of the
error matrix. Then $A$ is invertible. Its rows belong to the complex
span of the real vectors $M(\rho)$, proving their full real rank.
The input is Theorem 1.2 of W.~Banks,
\href{https://arxiv.org/abs/2302.07073}{\emph{Comparing zeros of distinct
Dirichlet $L$-functions}}. The conductor-one case is the classical
Landau formula in the same paper.\hfill$\square$

This yields a nonsingular integer $d$-by-$d$ multiplicity matrix in
every such height window, including when different characters share zeros.
In particular one can choose $\gg_q T/\log T$ pairwise disjoint
bases with ordinates in $[T,2T]$, by dividing it into intervals of
length $2K_q\log(2T)$. Under Dirichlet GRH these are bases of
multiplicity vectors attached to distinct heights. 
Banks also obtains bounds uniform in the conductor for a pair of
characters. The rank argument here keeps $q$ fixed. The entries of the
matrix record multiplicities. Its nonsingularity places no condition
on rational linear relations among the heights themselves.

\subsection{Isolating oscillations with positive weights}

For a nonzero ordinate $\gamma$ in the union of the spectra put
$\rho_\gamma=1/2+i\gamma$ and
\[
 c_{\gamma,a}=-\frac{1}{d\rho_\gamma}
               \sum_\chi \overline{\chi(a)} M(\rho_\gamma)_\chi,
 \qquad c_\gamma=(c_{\gamma,a})_{a\in G}.
\]
The finite Fourier transform and the rank theorem show that the
vectors $c_\gamma$ for positive $\gamma$ span $\mathbb C^d$.
Thus the real and imaginary parts of finitely many of them span $\mathbb R^d$.

On $v>1$, the explicit formula, as an identity of distributions, reads
\[
 D(e^v)=\sum_{\gamma\ne0}c_\gamma e^{i\gamma v}
          +O_q((1+v)e^{-v/2}).\tag{E}\label{eq:explicitvec}
\]
To see the error term, first use a primitive character. The main pole
gives $e_\chi e^v$, each nontrivial zero contributes $-e^{\rho v}/\rho$
with multiplicity, and the central contribution is canceled by
\eqref{eq:racevec}. The pole of the Perron kernel at zero and a possible
trivial zero there contribute $O_q(1+v)$ before normalization.
The other trivial-zero terms form a geometrically decreasing series
for $v>1$. Replacing a primitive character by the one modulo $q$
changes $\psi$ by $O_q(v+1)$, since only powers of primes dividing $q$
are affected. The choice at a jump is immaterial to integration.
A primary reference for the boundary-value formula and its convergence
for smooth averages is Kaczorowski and Ramar\'e, Theorem 1(d),(e),
\href{https://doi.org/10.4064/aa106-3-6}{\emph{Almost periodicity of some
error terms in prime number theory}}.

Take a smooth probability density $w$, positive on $(0,1)$ and supported
on $[0,1]$, and let $w_n$ be its $n$-fold convolution, supported on
$[0,n]$. Put $\widehat w(t)=\int w(u)e^{itu}\,du$. For each nonzero
real $t$, $|\widehat w(t)|<1$, and $\widehat w$ decreases rapidly.
For a fixed $\omega$, integration of (E) shows
\begin{align*}
 &\int_0^n w_n(u) D(e^{V+u})e^{-i\omega(V+u)}\,du \\
 &\qquad=\sum_{\gamma\ne0}
      c_\gamma e^{i(\gamma-\omega)V}\widehat w(\gamma-\omega)^n
       + O_q((1+V)e^{-V/2})\qquad(V\ge 2).
\end{align*}
The error is uniform in $n$ since $(1+v)e^{-v/2}$ decreases for $v\ge2$.
The sum converges absolutely for $n\ge1$ by rapid decay and the
classical $O_q(U\log(2U))$ zero count up to height $U$. The union of
the spectra is locally finite. Hence for each fixed $\omega$ the
supremum of $|\widehat w(\gamma-\omega)|$ over $\gamma\ne\omega$ is
strictly below $1$. All off-frequency terms tend uniformly in $V$ to
zero as $n\to\infty$. At $\omega=0$ the sum therefore tends to zero.
At a selected $\omega=\gamma$ it tends to $c_\gamma$.
For the arithmetic remainder take first a sufficiently large fixed
$n$, then $V$ sufficiently large.

Average $D(e^v)$ against the nonnegative densities proportional to
$w_n(v-V)(1\pm\cos(\gamma v))$ and
$w_n(v-V)(1\pm\sin(\gamma v))$. Their masses tend uniformly to $1$
because $|\widehat w(\gamma)|<1$. The averages approach, to any
prescribed accuracy, the vertices $\pm\Re c_\gamma$ and
$\pm\Im c_\gamma$ for the selected finite set of positive ordinates.
Here the opposite signs for sine merely interchange the two vertices.
The convex hull of these ideal vertices contains a ball about the
origin because their real span has full dimension. Concretely, if its
inradius is $r>0$, approximating every vertex within $r/4$ preserves
a ball of radius $3r/4$, by the support-function characterization of
convex inclusion.

Finally replace $D(e^v)$ in the averages by $D(\lfloor e^v\rfloor)$.
Between consecutive integers the numerator's arithmetic sum is
constant, and $\psi(x;q,a)\leq x\log(2x)$. The mean value theorem
bounds the change by $O_q((1+v)e^{-v/2})$, uniformly for large $v$.
The averages now belong to the convex hull of a \emph{finite} set of
integer samples. Set $V=\log(2X)$, so those integers lie in
$[X,2e^nX]$. Taking $C_q=2e^n$ and $r_q=r/2$ proves
\eqref{eq:ballrace}. Carath\'eodory's theorem gives the sample bound.
\hfill$\square$

\paragraph{Relation to the open prime race problem.}
A ball in the convex hull allows thresholds to be mixed. A prescribed
ranking of raw prime counts asks for all of its inequalities to hold
at one threshold. Even a circle has a convex hull containing points
that the circle never visits. For results on the single-threshold
problem under hypotheses on the zero ordinates, see G.~Martin and
N.~Ng, \href{https://doi.org/10.1090/tran/7996}{\emph{Inclusive prime
number races}}. The conditional almost-periodicity theory of Kaczorowski and Ramar\'e
(Theorem 1(e)) gives a second perspective. Their corrected vector has
Stepanov almost-periodic coordinates, zero mean after removing the
central terms, and the nonzero Fourier coefficients displayed in (E),
up to an overall sign. The rank theorem makes every nonzero real
projection nonconstant. Positive and negative values of a nonzero
mean-zero projection then recur on sets of positive measure in
intervals of bounded log-length: an interval on which it exceeds a
positive level on a positive-measure subset can be repeated using
relatively dense Stepanov almost-periods; the same applies to its
negative. The positive compact averages above give a direct proof
uniform across directions. This classical theory and Banks's
explicit formula account for the analytic inputs. We have not
established priority for the rank or convex-hull formulations.

\paragraph{An exact finite example.}
Take $q=4$. Its two central orders vanish. The nonprincipal central
value is a positive alternating series. For the principal character,
$\zeta(1/2)=\eta(1/2)/(1-\sqrt2)$ and the alternating eta series is
positive. The three integers $293,283,228$ lie in $[200,400]$.
For $v(n)=(\psi(n;4,1)-n/2,\psi(n;4,3)-n/2)$ the certified lower
bounds for the cyclic determinants are
\[
 \det(v(293),v(283))\ge5.625,\quad
 \det(v(283),v(228))\ge41.693,\quad
 \det(v(228),v(293))\ge14.050.
\]
All three are positive. The identity
$\det(v_2,v_3)v_1+\det(v_3,v_1)v_2+\det(v_1,v_2)v_3=0$
therefore gives strictly positive balancing weights. On replacing
$v(n)$ by $D(n)=v(n)/\sqrt n$, multiply the associated weight by
$\sqrt n$ and renormalize. The file
\texttt{checks/prime\_race\_triangle.py} verifies the three bounds using
integer sieving and rational upper and lower estimates for logarithms.
This example is independent of GRH.


\section{Further research questions}
\begin{enumerate}
\item Does every singular modulus of degree at least five require at least six terms? The full four-term result and the five-root certificates at 28 individual discriminants leave a concrete test case at $\Delta=-131$. More generally, can one certify $s(\Delta)$ from a combination of disks at different places?
\item Can the lower bound $s(\Delta)\gg\sqrt{h(\Delta)}$ be improved to a larger power of $h$? The proof balances the number of distinct residues against the size of a disk on which sparse powers are independent.
\item Does the ordinary logarithmic density in Suzuki's unsmoothed question exist at the specified threshold, and what is its value? The uniform bounds in Section~\ref{sec:suzuki-density} leave room for a small positive exceptional density.
\item For Suzuki's equation (29), what is the sign in the intermediate regime $Q\to\infty$ with $Q=o(x)$? For (30), can one replace the finite-logarithmic-measure exceptional set by an eventual pointwise statement at fixed $Q$?
\item Can the convex-hull radius and interval ratio in Section~\ref{sec:balancing} be made effective in terms of the modulus? A ranking at one sample would require stronger information than convex balance.
\end{enumerate}

\clearpage
\begin{thebibliography}{10}
\bibitem{upstream} L. Sernau and collaborators, \emph{GRH for Dirichlet $L$-functions, in Lean 4}, based on the method of L. Alp\"oge. \href{https://github.com/lsernau/grh-lean/tree/72fead102f663d5c7a407964a13dc0cf9d977207}{Pinned repository and build record}.
\bibitem{prior} ProveIt incoming report, \emph{Sparse CM equations, prime-generating constants, and additive representations}, November 25, 2026. \href{https://github.com/VladimirReshetnikov/ProveIt/blob/c59bde1a0f52eee7a7a8dc6b326072f9b71f76a6/docs/incoming/grh_cm_prime_applications.zip}{Archived manuscript and candidate list}.
\bibitem{avk} M. Avenda\~no and T. Krick, \emph{Sharp bounds for the number of roots of univariate fewnomials}, J. Number Theory 131 (2011), 1209--1228. \href{https://arxiv.org/abs/1008.4808}{Author manuscript}.
\bibitem{blp} Y. Bilu, F. Luca and A. Pizarro-Madariaga, \emph{Trinomials, singular moduli and Riffaut's conjecture}, Ann. Sc. Norm. Super. Pisa Cl. Sci. 23 (2022), 2003--2048. \href{https://arxiv.org/abs/2003.06547}{Author manuscript}.
\bibitem{cox} D. A. Cox, \emph{Primes of the Form $x^2+ny^2$}, 2nd ed., Wiley, 2013. See sections 8--11 for class fields and complex multiplication, and Theorem 7.24 for the order class number formula.
\bibitem{eoy} N. Elkies, K. Ono and T. Yang, \emph{Reduction of CM elliptic curves and modular function congruences}, Int. Math. Res. Not. 2005, 2695--2707. \href{https://arxiv.org/abs/math/0512350}{Author manuscript}.
\bibitem{lenstra} H. W. Lenstra, Jr., \emph{On the factorization of lacunary polynomials}, in \emph{Number Theory in Progress}, 1999, 277--292. \href{https://pub.math.leidenuniv.nl/~lenstrahw/PUBLICATIONS/1999b/art.pdf}{Author copy}.
\bibitem{silverman} J. H. Silverman, \emph{The Arithmetic of Elliptic Curves}, 2nd ed., Springer, 2009. Chapter V, \S4, gives the Hasse invariant test.
\bibitem{suzuki} M. Suzuki, \emph{On variants of Chebyshev's conjecture}, Ramanujan J. 68 (2025), art. 95, with \href{https://doi.org/10.1007/s11139-025-01289-y}{correction}. \href{https://arxiv.org/abs/2411.07436v3}{Corrected author manuscript}.
\bibitem{platt} D. J. Platt, \emph{Numerical computations concerning the GRH}, Math. Comp. 85 (2016), 3009--3027. \href{https://doi.org/10.1090/mcom/3077}{Paper}.
\bibitem{pt} D. Platt and T. Trudgian, \emph{The Riemann hypothesis is true up to $3\cdot10^{12}$}, Bull. Lond. Math. Soc. 53 (2021), 792--797. \href{https://doi.org/10.1112/blms.12460}{Paper}.
\bibitem{hsw} E. Hasanalizade, Q. Shen and P.-J. Wong, \emph{Counting zeros of the Riemann zeta function}. \href{https://arxiv.org/abs/2107.06506}{Author manuscript}.
\bibitem{banks} W. D. Banks, \emph{Comparing zeros of distinct Dirichlet $L$-functions}. \href{https://arxiv.org/abs/2302.07073}{Author manuscript}.
\bibitem{kr} J. Kaczorowski and O. Ramar\'e, \emph{Almost periodicity of some error terms in prime number theory}, Acta Arith. 106 (2003), 277--297. \href{https://doi.org/10.4064/aa106-3-6}{Paper}.
\end{thebibliography}
\end{document}
